\section*{Chapter \ref{ch:b1:chemical-biogenic-sediments} --- Carbonate, Siliceous and Evaporite Sediments}

\begin{solution}{exo:b1:chemical-biogenic-sediments:1}
Chloride $0.5503\times35.165 = \qty{19.35}{g}$; sodium $0.3066\times35.165 =
\qty{10.78}{g}$; sulfate \qty{2.711}{g}; magnesium \qty{1.284}{g}. Sodium and
chloride together make $0.5503 + 0.3066 = 0.857$, that is \qty{85.7}{\%} of the
salt.
\end{solution}

\begin{solution}{exo:b1:chemical-biogenic-sediments:2}
The fraction left is $1/\mathrm{CF}$: calcium carbonate $1/1.8 = 0.556$ left,
\qty{44.4}{\%} evaporated; gypsum $0.263$ left, \qty{73.7}{\%} evaporated;
halite $0.094$ left, \qty{90.6}{\%} evaporated; magnesium sulfate $0.014$ left,
\qty{98.6}{\%} evaporated.
\end{solution}

\begin{solution}{exo:b1:chemical-biogenic-sediments:3}
(a)~$\rho_s(1-\phi) = G/a = 8.4/0.0042 = \qty{2000}{kg.m^{-3}}$. (b)~With
aragonite, $1 - 2000/2930 = 0.32$, a porosity of \qty{32}{\%}; with calcite,
$1 - 2000/2710 = 0.26$. (c)~$a = 3.7/2000 = \qty{1.85e-3}{m}$, \qty{1.85}{mm} a
year.
\end{solution}

\begin{solution}{exo:b1:chemical-biogenic-sediments:4}
Down to \qty{4000}{m} the content falls only slowly (88 to 80); below
\qty{4000}{m} it falls steeply: the \omterm{def:b1:chemical-biogenic-sediments:ccd}{lysocline} lies at about \qty{4000}{m}. The
content falls through \qty{10}{\%} between \qty{4700}{m} (12) and \qty{4800}{m}
(6): $4700 + 100\times(12-10)/(12-6) = \qty{4733}{m}$, a CCD at about
\qty{4730}{m}, some \qty{700}{m} below the \omterm{def:b1:chemical-biogenic-sediments:ccd}{lysocline}.
\end{solution}

\begin{solution}{exo:b1:chemical-biogenic-sediments:5}
(a)~\ce{Ca^2+ + 2 HCO3- -> CaCO3 + CO2 + H2O}: one mole of \ce{CO2} per mole of
\ce{CaCO3}, so $44.01/100.09 = \qty{0.440}{kg}$ of \ce{CO2} per kilogram of
\ce{CaCO3}. (b)~$3.7\times0.440 = \qty{1.63}{kg}$ of \ce{CO2} per square metre and
per year. (c)~Of the two hydrogen carbonate ions used, only one carbon goes into
the \omterm{def:b1:rock-forming-minerals:mineral}{mineral}; the other leaves as carbon dioxide. Calcification takes alkalinity
out of the water and lowers its capacity to hold dissolved carbon; how much of
the \ce{CO2} reaches the air depends on the buffering of seawater, a matter of
the carbonate equilibria of the chemistry books.
\end{solution}

\begin{solution}{exo:b1:chemical-biogenic-sediments:6}
(a)~$t_c = (2000/350)^2 = \qty{32.7}{Myr}$; at \qty{25}{km} per million years,
$25\times32.7 = \qty{816}{km}$ from the axis. (b)~$t_c = (2600/350)^2 =
\qty{55.2}{Myr}$, \qty{1380}{km}. (c)~Floor \qty{60}{Myr} old has received red
clay for $60 - 32.7 = \qty{27.3}{Myr}$ in the first case, and for only
$60 - 55.2 = \qty{4.8}{Myr}$ in the second.
\end{solution}

\begin{solution}{exo:b1:chemical-biogenic-sediments:7}
(a)~Yes: $a = 4.2 > r = 3$, and the \omterm{def:b1:chemical-biogenic-sediments:reef}{reef} gains \qty{1.2}{mm} a year on the sea
until its top reaches the surface. (b)~$D$ grows by $10 - 4.2 = \qty{5.8}{mm}$ a
year; from 2 to \qty{50}{m} it takes $48/0.0058 = \num{8300}$ years. (c)~Because
its growth rate falls as the water deepens: less light reaches the algae that
live in the corals and feed them, so $a$ decreases and the \omterm{def:b1:chemical-biogenic-sediments:reef}{reef} falls behind
faster and faster.
\end{solution}

\begin{solution}{exo:b1:chemical-biogenic-sediments:8}
(a)~The reaction is \ce{CaSO4.2H2O -> CaSO4 + 2 H2O}. (b)~The molar volumes are $172.17/2.30 =
\qty{74.9}{cm^3}$ and $136.14/2.97 = \qty{45.8}{cm^3}$, so the solid loses
$1 - 45.8/74.9 = \qty{38.8}{\%}$ of its volume. (c)~$2\times18.02/172.17 = 0.209$:
\qty{20.9}{\%} of the mass leaves as water. Released at depth, this water
raises the pressure in the pores of the rocks around; under a seal it can make
the rocks weaker and help faults slip, and it can dissolve and carry other
salts.
\end{solution}

\begin{solution}{exo:b1:chemical-biogenic-sediments:9}
(a)~$9.2/255 = \qty{3.6}{\%}$ preserved, \qty{96.4}{\%} dissolved. (b)~Where
production is very high, a few per cent of it is still a large flux, and where
little clay or carbonate dilutes it: under the upwelling waters of the Southern
Ocean, the North Pacific and the equatorial Pacific, and below the CCD.
(c)~Thin, delicate frustules dissolve first and robust ones survive: the
sediment keeps a selection of the sturdiest species, not the community that
lived in the water.
\end{solution}

\begin{solution}{exo:b1:chemical-biogenic-sediments:10}
(a)~$EA = 1.5\times\qty{1e10}{m^2} = \qty{1.5e10}{m^3}$ a year. (b)~With
$Q_{\mathrm{in}} = \qty{1.8e10}{m^3}$, $Q_{\mathrm{out}} = \qty{0.3e10}{m^3}$ and
$k = 6$: past the gypsum threshold (3.8) but not that of halite (10.6), so
calcium carbonate and gypsum precipitate. With $Q_{\mathrm{in}} =
\qty{1.575e10}{m^3}$, $Q_{\mathrm{out}} = \qty{0.075e10}{m^3}$ and $k = 21$, for
an ion that stays in solution such as magnesium: the brine passes the halite
threshold, so halite precipitates too, but not the magnesium salts (70).
(c)~Halite requires $Q_{\mathrm{in}}/(Q_{\mathrm{in}} - EA) \ge 10.6$, that is
$Q_{\mathrm{in}} \le 10.6\,EA/9.6 = 1.104\,EA = \qty{1.66e10}{m^3}$ a year: the
inflow may exceed the evaporation by at most about a tenth.
\end{solution}

\begin{solution}{exo:b1:chemical-biogenic-sediments:11}
(a)~Seawater holds \qty{67.3}{mg} of bromide per kilogram; at 10.6~times, the
brine holds $67.3\times10.6 = \qty{713}{mg.kg^{-1}}$ and the first halite
$0.032\times713 = \qty{22.8}{mg.kg^{-1}}$. (b)~At 30~times, $67.3\times30 =
\qty{2019}{mg.kg^{-1}}$ in the brine and \qty{64.6}{mg.kg^{-1}} in the halite.
(c)~The halite takes only 3.2~hundredths of the concentration of the brine, so
nearly all of the bromide stays in solution and keeps concentrating as the water
evaporates. (d)~The rise records the concentration of the brine; the sudden fall
records its dilution, by a new inflow of seawater into the basin, or of fresh
water, before the halite went on growing.
\end{solution}

\begin{solution}{exo:b1:chemical-biogenic-sediments:12}
(a)~$t_c = \qty{32.65}{Myr}$, so $0.02\times\num{32650} = \qty{653}{m}$ of
ooze (\num{32650} thousand-year intervals at \qty{2}{cm}) and
$0.001\times\num{37350} = \qty{37}{m}$ of red clay. (b)~$t_c = (2500/350)^2 =
\qty{51.0}{Myr}$: \qty{1020}{m} of ooze and \qty{19}{m} of red clay. A deeper CCD
gives a much thicker ooze, and the red clay of the deep Atlantic floor is thinner
than that of the Pacific of the same age. (c)~The CCD may have been deeper for a
while, so that the site rose above it again; or the plate carried the site under
a belt of high productivity, such as the equator, where more carbonate rained
down; or the layer was brought by a slide or a \omterm{def:b1:clastic-sediments:turbidite}{turbidity current} from a higher
part of the sea floor.
\end{solution}

\begin{solution}{pb:b1:chemical-biogenic-sediments:1}
\textbf{1.} $1030\times1000 = \qty{1.03e6}{kg}$ of seawater, holding
$1.03\times10^6\times0.035165 = \qty{36200}{kg}$ of salt.

\textbf{2.} Gypsum $1.03\times10^6\times1.581\times10^{-3} = \qty{1628}{kg}$;
halite $1.03\times10^6\times0.02741 = \qty{28230}{kg}$.

\textbf{3.} Calcium carbonate $113.3/2710 = \qty{0.042}{m}$; gypsum $1628/2300 =
\qty{0.71}{m}$; halite $28230/2170 = \qty{13.0}{m}$.

\textbf{4.} The anhydrite has the mass $1628\times136.14/172.17 =
\qty{1288}{kg}$ and the thickness $1288/2970 = \qty{0.43}{m}$.

\textbf{5.} Dry, the salt is $0.110 + 1.250 + 27.41 + 6.34 = \qty{35.11}{g}$ per
kilogram, of which halite is $27.41/35.11 = \qty{78}{\%}$: sodium and chloride
make most of the dissolved salt, and halite takes nearly all of them.

\textbf{6.} From the bottom up: calcium carbonate, gypsum, halite, the bittern
salts. Each \omterm{def:b1:rock-forming-minerals:mineral}{mineral} starts to precipitate when the brine reaches its
saturation, the least soluble first, and settles on what was laid down before.

\textbf{7.} The depth of brine is $1000/\mathrm{CF}$: \qty{556}{m} for calcium
carbonate, \qty{263}{m} for gypsum, \qty{94}{m} for halite.

\textbf{8.} $1 - 1/10.6 = \qty{90.6}{\%}$.

\textbf{9.} $1000/70 = \qty{14}{m}$ for magnesium sulfate, $1000/90 = \qty{11}{m}$
for the potassium salts.

\textbf{10.} About \qty{1.10}{g.cm^{-3}} when calcium carbonate forms, 1.13 with
gypsum, 1.22 with halite and more than 1.30 with the bitterns, against 1.03 for
seawater. Seawater flowing in floats on the dense brine, which stays at the
bottom; in a barred basin the brine can flow back to the ocean below the
inflow.

\textbf{11.} It is only \qty{4}{cm} thick, against \qty{13}{m} of halite; and it
forms on the margins, where the inflowing water first concentrates, more than in
the centre, so it is mixed with the base of the gypsum or replaced by it.

\textbf{12.} The halite takes only a small part of the bromide of the brine
(exercise~\ref{exo:b1:chemical-biogenic-sediments:11}), so the brine becomes
richer in bromide as the halite grows, and each new layer of halite takes up
more.

\textbf{13.} $1500/13.01 = 115.3$: 116 fillings.

\textbf{14.} $116\times0.434 = \qty{50}{m}$ of anhydrite, not as one bed but as
116 layers, each a little under half a metre thick, at the base of the halite of
its filling: a cyclic succession of anhydrite and halite.

\textbf{15.} Per filling, $1.03\times10^6\times6.34\times10^{-3}/2000 =
\qty{3.3}{m}$, so \qty{379}{m} for 116 fillings. They are the most soluble
salts: the next inflow of seawater redissolves them, and a basin whose brine
flows out before the end never reaches them.

\textbf{16.} At each refilling the brine is diluted: the bromide content drops
sharply at the base of the new halite and rises again upward, a saw-tooth
profile with one tooth per filling.

\textbf{17.} $2000/13.01 = 154$ fillings.

\textbf{18.} $5000\times365.25\times86400 = \qty{1.578e11}{m^3}$ a year, spread
over $\qty{1.0e11}{m^2}$: $E = \qty{1.58}{m}$ a year.

\textbf{19.} $1.578\times10^{11}\times1030\times0.035165 = \qty{5.72e12}{kg}$ of
salt a year.

\textbf{20.} $1.578\times10^{11}\times1030\times0.02741/2170 =
\qty{2.05e9}{m^3}$ of halite a year, that is \qty{0.0205}{m}, about \qty{2}{cm}
a year over the floor.

\textbf{21.} $1500/0.0205 = \num{73000}$ years.

\textbf{22.} The level falls at \qty{1.58}{m} a year; \qty{1000}{m} of water is
gone in $1000/1.578 = 634$ years.

\textbf{23.} Halite appears when \qty{90.6}{\%} of the water has gone, after
$905.7/1.578 = 574$ years.

\textbf{24.} The \qty{1.5}{km} of halite took the salt of 116 fillings of the
basin (115.3, to be rounded up), whether they came one after the other as the
basin dried and refilled, or continuously over the sill, as in
question~21, at the rate of one basin of seawater every
$1000/1.578 = 634$ years.
\end{solution}
